A manipulação de bits é uma técnica fundamental na programação em C que permite operar 
diretamente nos bits individuais de variáveis. Essas operações são extremamente eficientes 
e são amplamente utilizadas em sistemas embarcados, programação de baixo nível, criptografia, 
compressão de dados e otimização de performance.

Já que é possível complicar as coisas facilmente quando se trata de operações envolvendo manipulação 
de valores em binários diretamente, vou tentar manter a explicação simples e direta com 
alguns exemplos sobre o assunto:

\subsection{Operadores Bitwise}

C fornece seis operadores para manipulação de bits, que operam em cada bit individualmente:

\begin{itemize}
    \item \texttt{\~} (NOT) - Complemento de bits
    \item \texttt{\&} (AND) - E bit a bit
    \item \texttt{|} (OR) - OU bit a bit
    \item \texttt{\^} (XOR) - OU exclusivo bit a bit
    \item \texttt{<<} (Left Shift) - Deslocamento à esquerda
    \item \texttt{>>} (Right Shift) - Deslocamento à direita
\end{itemize}

\newpage
\begin{excode} #include <stdio.h>
    ...
    int main(void)
    {

        // a = 5 (00000101 em binário)
        // b = 9 (00001001 em binário)
        unsigned int a = 5, b = 9;

        // O resultado é 00000001
        printf("a&b = %u\n", a & b);

        // O resultado é 00001101
        printf("a|b = %u\n", a | b);

        // O resultado é 00001100
        printf("a^b = %u\n", a ^ b);

        // O resultado é 11111111111111111111111111111010
        // (assumindo um inteiro de 32 bits)
        printf("~a = %u\n", a = ~a);

        // O resultado é 00010010
        printf("b<<1 = %u\n", b << 1);

        // O resultado é 00000100
        printf("b>>1 = %u\n", b >> 1);
        return 0;
        }
\end{excode}

\subsection{Truques úteis}

\subsubsection{Definir/inverter/limpar um bit}

Usando deslocamento bit-a-bit é possível definir/inverter/limpar bits facilmente.

\begin{excode}
    ...
    // Definir um número com apenas o x-ésimo bit como 1.
    int a = 1 << x;
    // a = 16
    int a = 1 << 4;

    // Definir um número com todos os bits como 1, 
    // exceto o x-ésimo bit.
    int a = ~(1 << x);
    // a = 4294967279
    int a = ~(1 << x);
    ...
\end{excode}

\subsubsection{Verificando um bit específico}

O valor do x-ésimo bit pode ser verificado deslocando o número x posições para a 
direita, em seguida, realizamos uma operação AND (\&) com 1.

\begin{excode}
    ...
    bool checar_bit(unsigned int numero, int index)
    {
        return (number >> x) & 1;
    }
    ...
\end{excode}

\subsubsection{Verificando se o número é ímpar ou par}

Usando uma simples verificação de AND (\&) é possível verificar se um número é ímpar 
ou par sem precisar de nenhum outro tipo de verificação mais complexa.

\begin{excode}
    ...
    bool divisivel_por_dois(int n)
    {
        return (n & 1);
    }
    ...
\end{excode}

\begin{postscript}
    A vantagem de usar o operador de AND para verificar se um número é ímpar ou par ao invés de 
    \texttt{return (n \% 2) == 0;} é que o compilador da linguagem C consegue otimizar o código muito mais, 
    o que pode ser muito útil em sistemas críticos que necessitem do máximo de desempenho possível.

    A nível de curiosidade, essa é a diferença entre ambos os códigos quando compilados pelo 
    \texttt{x86-64 gcc 15.2} usando a flag \texttt{-O3}.

    Usando o método do operador AND:
    \begin{exterminal}
        ...
        div_by_2_bitwise:
        mov     esi, edi
        test    dil, 1
        je      .L2
        mov     edi, OFFSET FLAT:.LC0
        xor     eax, eax
        jmp     printf
        ...
    \end{exterminal}

    Usando o método do operador AND:
    \begin{exterminal}
        ...
        div_by_2_percent:
        mov     edx, edi
        mov     esi, edi
        shr     edx, 31
        lea     eax, [rdi+rdx]
        and     eax, 1
        sub     eax, edx
        cmp     eax, 1
        je      .L6
        mov     edi, OFFSET FLAT:.LC1
        xor     eax, eax
        jmp     printf
        ...
    \end{exterminal}

    Mesmo sem conhecimentos em Assembly é possível notar uma diferença gritante na 
    quantidade de instruções para que a função seja executada corretamente.
\end{postscript}
